\section{Requirement-Aware Physical-Plan Allocator}
\label{sec:optimizer}

% OP-P0 — Repositioning framework (P2 path b): the general enummeration is the
% real mechanism; on the chain topology the allocator is deterministic.
Section~\ref{sec:algebra} gives requirements importance, type, and dependency.
This section describes how slotrag turns a logical evidence program into a
physical plan under a matched budget, and how the physical-plan of the chain
workload we study reduces to a deterministic, importance-weighted allocation.

We adopt an explicit \emph{prediction/decision} separation: a property estimator
supplies estimates of optimizer-visible properties (yield, coverage, cost), and
a physical-plan allocator turns them into an execution order, a retrieval-call
allocation, and a physical implementation per requirement. On the well-defined
chains that dominate our workloads, requirement importance is
construction-derived and exact ($\tau = 2d - 1$), so the allocation is
deterministic and auditable rather than a searched optimum. On general
(non-chain) graphs the allocator falls back to a bounded enumeration over legal
orders, physical implementations, and budget allocations; because the evaluated
workloads are chain-topology, the enumeration is a \emph{mechanism}, not a claim
of general plan search --- the branching re-optimization that would give
enumeration non-trivial freedom is future work (Section~\ref{sec:execution}).
We state this boundary so the reader does not over-read general plan-searching
into a system that today runs chains.

% OP-P2 — Estimator architecture
The optimizer makes decisions from estimates of optimizer-visible properties:
expected evidence yield, marginal requirement coverage, selectivity, failure
probability, and cost. Our design deliberately separates \emph{prediction} from
\emph{decision}: the estimator supplies property estimates; an explicit
optimizer consumes them. For the well-defined chains that dominate our
workload, we derive the requirement importance estimate in closed form. The
chain law states that a requirement at dependency depth $d$ carries importance
$\tau = 2d - 1$: each requirement is worth the depth of the chain it anchors,
and the downstream requirement is strictly more important than the upstream
one. This law was discovered by a counterfactual scan over compiled plans and
holds deterministically on well-defined chains (Spearman correlation 1). It is
a calibrated, construction-derived estimator rather than a learned one; we do
not claim a learned estimator adds value on this subdomain. On
non-well-defined chains the closed form is a proxy artifact and is outside the
claim.

% OP-P3 — Calibration / training data
Because the importance estimator is derived by construction, it requires no
training data and cannot leak sealed test information. This is the honest
fallback the design intends: when a property estimator admits no calibration
from development executions, the system uses a deterministic, construction-level
estimator rather than silently substituting a different learned model. All
evaluation in this paper respects the development / validation / sealed-test
registry; no sealed instance participates in any threshold tuning or property
calibration.

% OP-P4 — Plan enumeration
The optimizer enumerates candidate physical plans by combining three factors:
legal execution orders (dependency-respecting topological orders of the
requirement graph), compatible physical implementations (hybrid vs.\ bm25 per
requirement), and budget allocations (a distribution of the retrieval-call
budget across requirements). Enumeration is bounded: dependency orders are
capped at a deterministic subset ($256$) of topological orders, and each
candidate is deduplicated by a canonical key. For a strict chain, the
dependency graph admits exactly one legal order, so enumeration reduces to
budget allocation; for branching graphs, multiple orders become legal and the
allocation can differ across them.

% OP-P5 — Dominance / Pareto pruning
Candidates are pruned by dominance. A candidate plan $A$ dominates $B$ if $A$
has at least as much expected requirement utility and no more expected cost on
every constrained resource dimension. Dominance pruning is necessary because a
cheaper plan does not automatically dominate a more expensive one: a plan that
spends more calls on a high-importance downstream requirement can satisfy a
materially different evidence state. Pruning removes only candidates that are
Pareto-dominated, and the number of pruned candidates is recorded in the
optimizer telemetry so the search itself is auditable.

% OP-P6 — Utility-aware selection
Plan selection is requirement-aware. The objective is

\vspace{-0.1em}
\begin{equation}
\label{eq:objective}
U(\pi) = \sum_{s \in \pi} w_s \cdot \Bigl(1 - e^{-\,c_s / \beta_s}\Bigr) - \lambda(\pi),
\end{equation}
\vspace{-0.1em}

\noindent
where $w_s$ is the importance of requirement $s$ ($\tau = 2d-1$ on well-defined
chains), $c_s$ is the number of retrieval calls allocated to $s$, $\beta_s$
scales with the requirement's estimated cardinality (rarer evidence needs more
calls), and $\lambda(\pi)$ is the estimated resource cost of the plan. The
saturating marginal means utility is credited only for \emph{expected progress}
on a requirement: marginal value diminishes as calls accumulate, and an
already-satisfied requirement contributes no further utility. This discourages
repeated retrieval of already-satisfied facts even when such retrieval would
rank highly on relevance alone---the failure mode of pure cost- or
relevance-driven selection.

% OP-P7 — Complexity and algorithm
\begin{figure}[t]
\centering
\caption{Bounded physical-plan search (requirement-aware optimizer).}
\label{alg:search}
\small
\begin{tabular}{@{}l@{}l}
$C \leftarrow \emptyset$ & candidate plans \\
\textbf{for} each legal topo order $\pi$ of dependency edges $E$: & \\
\quad $\mathit{alloc} \leftarrow \textsc{Allocate}(\pi,\ \tau,\ B)$ & importance-weighted greedy \\
\quad \textbf{for} each physical assignment $\sigma$ per slot: & \\
\qquad $C \leftarrow C \cup \{(\pi,\ \sigma,\ \mathit{alloc})\}$ & \\
$\mathit{best} \leftarrow \arg\max_{c\in C} \mathrm{Utility}(c)$ & Pareto-dominance, tie: fewer calls \\
\textbf{return} $\mathit{best}$ &
\end{tabular}
\end{figure}

Algorithm~\ref{alg:search} presents the bounded search. The dominant cost
factors are the number of legal orders and the number of physical
implementations per requirement; both are capped (orders $\le 256$,
implementations $\le 2$ per requirement). The allocation step is a deterministic
greedy pass over requirements in a fixed order. In the worst case the search is
exponential in the number of requirements before the cap, which is why the cap
is a declared search-space bound rather than a silent truncation: the
requirement-aware selection is robust to the subset because the dominant
candidate is always the dependency-ordered plan. The realized search overhead
is reported as a systems cost in Section~\ref{sec:mechanism}.

\paragraph{Complexity bound.}
On a strict chain the dependency graph admits exactly one legal execution order
($\pi$ is fixed), so the enumeration factor collapses and the allocator reduces
to the deterministic importance-weighted greedy over a fixed order
(\S OP-P4); search overhead is then $O(|S|)$ for allocation plus $O(2^1)$
physical-assignment checks, which is dominated by a single retrieval forward
pass. On a branching graph the number of legal orders $O$ and per-requirement
implementations $I$ are both bounded ahead of time ($O \le 256$, $I \le 2$),
giving a declared worst case of $O(256 \cdot 2^{|S|} \cdot |S|)$ candidate
plans; the Pareto-pruning step (\S OP-P5) removes dominated candidates during
enumeration, and the number of pruned candidates is recorded in optimizer
telemetry. Because the evaluated workloads are chain-topology, the realized
cost sits at the $O(|S|)$ end of this bound rather than the declared worst case.
